% Chapter 18: Streamlines and Euler's Equations for Fluids
% ch18_streamlines.tex

\chapter{Streamlines and Euler's Equations for Fluids}

\begin{historical}
In 1755, Euler published \emph{Lezioni di statica dei fluidi}, deriving the equations of motion for an ideal (inviscid) fluid from first principles. He showed that for a fluid with no viscosity, the acceleration of a fluid particle equals the sum of pressure and body forces per unit mass. These \textbf{Euler equations} are the foundation of inviscid fluid dynamics, and the concept of a \textbf{streamline}---a curve everywhere tangent to the velocity field---was Euler's key insight for visualizing and computing fluid flow.

Euler's approach was remarkable for its time: he applied Newton's second law directly to a fluid element, treating the fluid as a continuum rather than a collection of particles. This continuum approach, now called \textbf{Eulerian mechanics}, contrasts with the \textbf{Lagrangian} approach that tracks individual fluid particles. Both perspectives remain fundamental in modern fluid dynamics.
\end{historical}

\section{Prerequisites}

To work through this chapter, you need:
\begin{itemize}
    \item \textbf{Mathematics:} Vector calculus basics (gradient, divergence), ordinary differential equations, and the concept of a velocity field $\vec{v}(x,y)$.
    \item \textbf{Programming:} Python functions as first-class objects, list comprehensions, and numerical integration (Euler's method or RK4).
\end{itemize}

\section{Derivation of Euler's Equations}

Euler's equations follow directly from Newton's second law applied to a fluid element.
Consider a small fluid element of volume $dV$ and density $\rho$. The forces acting on it are:
\begin{itemize}
    \item \textbf{Pressure forces:} $-\nabla p \, dV$ (acting inward, opposite to pressure gradient)
    \item \textbf{Body forces:} $\rho \vec{g} \, dV$ (e.g., gravity)
\end{itemize}

The acceleration of the fluid element is given by the material derivative:
\[
\frac{D\vec{v}}{Dt} = \frac{\partial \vec{v}}{\partial t} + (\vec{v} \cdot \nabla)\vec{v}
\]

For steady flow, $\partial \vec{v}/\partial t = 0$, so the acceleration is purely convective.
Newton's second law gives:
\[
\rho \frac{D\vec{v}}{Dt} = -\nabla p + \rho \vec{g}
\]

This is the \textbf{Euler equation} for inviscid flow. For incompressible flow ($\rho = \text{const}$),
we can write it in the more familiar form.

\subsection{Component Form}

In two dimensions, with velocity components $u(x,y)$ and $v(x,y)$, the equations become:
\[
\rho \left( u \frac{\partial u}{\partial x} + v \frac{\partial u}{\partial y} \right) = -\frac{\partial p}{\partial x} + \rho g_x
\]
\[
\rho \left( u \frac{\partial v}{\partial x} + v \frac{\partial v}{\partial y} \right) = -\frac{\partial p}{\partial y} + \rho g_y
\]
where $u,v$ are velocity components, $p$ is pressure, $\rho$ is density, and $g_x, g_y$ are body forces.

In vector form: $\rho (\vec{v} \cdot \nabla) \vec{v} = -\nabla p + \rho \vec{g}$.

\begin{lstlisting}[language=Python, caption=Euler's steady-flow equations]
from euler import euler_equations_steady

# Given velocity field (u(x,y), v(x,y))
def vx(x, y): return 1.0 - x*x + y*y
def vy(x, y): return -2*x*y

# Compute pressure gradients
dpdx, dpdy = euler_equations_steady(vx, vy, rho=1.0)

# Evaluate at a point
p_grad_x = dpdx(0.5, 0.3)
p_grad_y = dpdy(0.5, 0.3)
\end{lstlisting}

\section{Streamlines}

A \textbf{streamline} is a curve $(x(s), y(s))$ everywhere tangent to the velocity field:
\[
\frac{dx}{ds} = u(x,y), \quad \frac{dy}{ds} = v(x,y)
\]
This is an ODE system. Integrating numerically from a starting point $(x_0, y_0)$ traces the streamline.

For steady flow, streamlines coincide with \textbf{pathlines} (trajectories of individual particles)
and \textbf{streaklines} (loci of particles passing through a fixed point). In unsteady flow,
these three concepts differ.

\subsection{Numerical Integration}

We use the fourth-order Runge-Kutta method (RK4) for accurate streamline computation:
\[
x_{n+1} = x_n + \frac{ds}{6}(k_1 + 2k_2 + 2k_3 + k_4)
\]
where $k_1 = u(x_n, y_n)$, $k_2 = u(x_n + \tfrac{ds}{2}k_1, y_n + \tfrac{ds}{2}k_2)$, etc.

\begin{lstlisting}[language=Python, caption=Computing a streamline]
from euler import streamline_integrate, uniform_flow

vx, vy = uniform_flow(vx_val=1.0, vy_val=0.5)

# Integrate from (0, 0)
pts = streamline_integrate(vx, vy, 0.0, 0.0,
                           ds=0.01, n_steps=200)

print(f"Streamline: {len(pts)} points")
print(f"End point: ({pts[-1][0]:.2f}, {pts[-1][1]:.2f})")
\end{lstlisting}

\section{Basic Flow Patterns}

Euler identified several elementary flows that can be superposed to model complex situations:

\textbf{Uniform flow:} $\vec{v} = (U, 0)$. Straight, parallel streamlines.

\textbf{Source/Sink:} $\vec{v} = \frac{m}{2\pi r} \hat{r}$. Radial flow outward (source) or inward (sink).

\textbf{Vortex:} $\vec{v} = \frac{\Gamma}{2\pi r} \hat{\theta}$. Circular streamlines around a center.

\textbf{Doublet:} Limit of a source-sink pair as separation $\to 0$. Used to model flow past cylinders.

Each of these flows is \textbf{irrotational} (zero vorticity) except at singularities,
meaning $\nabla \times \vec{v} = 0$. This allows us to define a \textbf{velocity potential}
$\phi$ such that $\vec{v} = \nabla \phi$.

\subsection{Stream Function}

For 2D incompressible flow, we can also define a \textbf{stream function} $\psi$ such that:
\[
u = \frac{\partial \psi}{\partial y}, \qquad v = -\frac{\partial \psi}{\partial x}
\]
The stream function automatically satisfies the incompressibility condition $\nabla \cdot \vec{v} = 0$.
Streamlines are contours of constant $\psi$.

\begin{lstlisting}[language=Python, caption=Superposing flow patterns]
from euler import (uniform_flow, source_flow,
                   vortex_flow, flow_past_cylinder,
                   streamline_integrate)

# Uniform flow + source
def vx_combo(x, y):
    u, _ = uniform_flow(1.0)
    s, _ = source_flow(strength=0.5)
    return u(x, y) + s(x, y)

def vy_combo(x, y):
    _, v = uniform_flow(1.0)
    _, sv = source_flow(strength=0.5)
    return v(x, y) + sv(x, y)

pts = streamline_integrate(vx_combo, vy_combo,
                           0.0, 1.0, ds=0.02, n_steps=300)
\end{lstlisting}

\section{Flow Past a Cylinder}

Superposing uniform flow with a doublet produces the classic potential flow around a circular cylinder. The stream function is:
\[
\psi = U \left(r - \frac{a^2}{r}\right) \sin\theta
\]
where $a$ is the cylinder radius and $U$ is the freestream velocity.

On the cylinder surface ($r = a$), we have $\psi = 0$, confirming that the cylinder is a streamline.
The velocity on the surface is purely tangential:
\[
v_\theta = -\frac{\partial \psi}{\partial r}\bigg|_{r=a} = -2U\sin\theta
\]

\begin{lstlisting}[language=Python, caption=Flow past a cylinder]
from euler import flow_past_cylinder, streamline_integrate

vx, vy = flow_past_cylinder(radius=1.0, freestream=1.0)

# Streamlines above and below the cylinder
for y0 in [0.5, 1.0, 1.5, 2.0]:
    pts = streamline_integrate(vx, vy, -3.0, y0,
                               ds=0.03, n_steps=200)
    print(f"y0={y0}: {len(pts)} points")
\end{lstlisting}

This produces \textbf{d'Alembert's paradox}: the drag force on a cylinder in potential flow is exactly zero.
Real fluids have viscosity, which creates boundary layers and wake formation.

\subsection{Stagnation Points}

Stagnation points occur where the velocity vanishes. For flow past a cylinder, the stagnation
points are at $(x,y) = (\pm a, 0)$. At these points, the pressure reaches its maximum value
(ram pressure).

\begin{lstlisting}[language=Python, caption=Finding stagnation points]
from euler import flow_past_cylinder, stagnation_points

vx, vy = flow_past_cylinder(radius=1.0, freestream=1.0)
stag_pts = stagnation_points(vx, vy, x_range=(-3, 3), y_range=(-3, 3))
print(f"Stagnation points: {stag_pts}")
\end{lstlisting}

\section{Bernoulli's Principle Along Streamlines}

For steady, inviscid, incompressible flow, Bernoulli's equation holds along any streamline:
\[
p + \frac{1}{2}\rho|\vec{v}|^2 = \text{constant}
\]
Higher velocity implies lower pressure. This explains lift on airfoils, the Venturi effect,
and why streamlines crowd together where velocity increases.

\subsection{Derivation from Euler's Equations}

Taking the dot product of Euler's equation with $\vec{v}$ and using the identity:
\[
(\vec{v} \cdot \nabla)\vec{v} = \nabla\left(\frac{|\vec{v}|^2}{2}\right) - \vec{v} \times (\nabla \times \vec{v})
\]

For irrotational flow ($\nabla \times \vec{v} = 0$), this simplifies to:
\[
\nabla\left(p + \frac{1}{2}\rho|\vec{v}|^2\right) = 0
\]
showing that $p + \frac{1}{2}\rho|\vec{v}|^2$ is constant everywhere (not just along streamlines)
for irrotational flow.

\begin{lstlisting}[language=Python, caption=Computing pressure from Bernoulli]
from euler import bernoulli_pressure, flow_past_cylinder

vx, vy = flow_past_cylinder(radius=1.0)
p = bernoulli_pressure(vx, vy, rho=1.2,
                       p_inf=101325.0, v_inf=10.0)

# Pressure at a point
p_at_point = p(1.5, 0.5)
print(f"Pressure at (1.5, 0.5): {p_at_point:.1f} Pa")
\end{lstlisting}

\section{Vorticity and Circulation}

The \textbf{vorticity} $\omega = \nabla \times \vec{v}$ measures local rotation of fluid elements.
For the elementary flows above (except the vortex singularity), vorticity is zero.

\textbf{Circulation} around a closed curve $C$ is defined as:
\[
\Gamma = \oint_C \vec{v} \cdot d\vec{r}
\]
By Stokes' theorem, $\Gamma = \iint_S \omega \, dA$ where $S$ is any surface bounded by $C$.
For irrotational flow, circulation around any closed curve is zero (unless the curve encloses a vortex).

\section{Modern Relevance}

Streamline theory remains essential in aerodynamics (airfoil design), hydrodynamics (ship hulls),
meteorology (weather map streamlines), and biomedical engineering (blood flow modeling).
Modern CFD solvers still solve the Euler equations as a first step before adding viscous terms.
The concept of streamlines is fundamental to every visualization of fluid flow, from wind tunnel
smoke trails to satellite weather maps.

\section{Frequently Asked Questions}

\begin{itemize}
\item \textbf{Q: What is the difference between a streamline, a pathline, and a streakline?} \\
For steady flow, all three coincide. A \textbf{streamline} is tangent to the instantaneous velocity field. A \textbf{pathline} traces the trajectory of a single fluid particle over time. A \textbf{streakline} connects all particles that have passed through a fixed point. Only in unsteady flow do they differ.

\item \textbf{Q: Why does potential flow predict zero drag on a cylinder?} \\
This is d'Alembert's paradox. The Euler equations assume zero viscosity, so there is no boundary layer separation and no wake. The pressure distribution is symmetric fore-and-aft, yielding zero net drag. Real fluids have viscosity, which breaks this symmetry.

\item \textbf{Q: How do you compute a streamline numerically?} \\
Treat the streamline ODE $dx/ds = u(x,y), dy/ds = v(x,y)$ as a system and integrate with any ODE solver (Euler's method, RK4). The step size $ds$ controls accuracy; smaller steps give smoother curves but require more computation.

\item \textbf{Q: What is the difference between a source and a vortex?} \\
A source has radial flow outward (or inward for a sink): $\vec{v} \propto \hat{r}/r$. A vortex has tangential flow: $\vec{v} \propto \hat{\theta}/r$. Both are irrotational except at the singularity. Sources model fluid injection; vortices model rotating flows.

\item \textbf{Q: Can streamlines cross?} \\
No. At any point in the flow, the velocity has a unique direction. If two streamlines crossed, a particle at the intersection would have two different velocities simultaneously, which is impossible. Streamlines can only meet at stagnation points where velocity is zero.

\item \textbf{Q: How does Bernoulli's equation relate to conservation of energy?} \\
Bernoulli's equation is the mechanical energy balance for a fluid particle. The term $p/\rho$ represents flow work (pressure energy), $\frac{1}{2}|\vec{v}|^2$ is kinetic energy per unit mass, and $gz$ is potential energy per unit mass. Their sum is constant along a streamline in inviscid, steady flow.

\item \textbf{Q: What is the significance of irrotational flow?} \\
Irrotational flow ($\nabla \times \vec{v} = 0$) allows the introduction of a velocity potential $\phi$ such that $\vec{v} = \nabla \phi$. This reduces the problem to solving Laplace's equation $\nabla^2 \phi = 0$, which is much simpler than the full Euler equations.
\end{itemize}

\begin{exercise}
Compute streamlines for a source of strength $m=1$ placed in a uniform flow of speed $U=1$.
Find the stagnation point where the velocity vanishes. Show that the streamline through the
stagnation point forms a semi-infinite body shape.
\end{exercise}

\begin{exercise}
Show that the stream function $\psi = U(r - a^2/r)\sin\theta$ for flow past a cylinder yields
zero normal velocity at $r=a$. Compute the tangential velocity on the cylinder surface.
\end{exercise}

\begin{exercise}
Use Bernoulli's equation to compute the pressure coefficient $C_p = (p - p_\infty)/(0.5\rho U^2)$
at the surface of a cylinder in potential flow. Where is $C_p$ minimum? What is the drag
force predicted by this model?
\end{exercise}

\begin{exercise}
Compute the circulation around a circle of radius $2a$ centered on a vortex of strength $\Gamma$
at the origin. Verify that the circulation equals $\Gamma$.
\end{exercise}

\begin{exercise}
Show that the superposition of a uniform flow $U$ and a doublet of strength $\kappa = 2\pi U a^2$
produces flow past a cylinder of radius $a$. Derive the stream function and velocity field.
\end{exercise}

\begin{exercise}
For the velocity field $u = x^2 - y^2$, $v = -2xy$, compute the vorticity and show that
the flow is irrotational. Find the velocity potential $\phi$ and stream function $\psi$.
\end{exercise}

\begin{exercise}
Compute the stagnation points for flow past a cylinder with circulation (the \textbf{Magnus effect}).
How does the circulation affect the lift force?
\end{exercise}

\begin{exercise}
Derive the velocity field for a source-sink pair (dipole) and show that it can be obtained
as the limit of a source and sink approaching each other.
\end{exercise}

\begin{exercise}
Use the stream function to compute the volume flow rate between two streamlines.
Show that the flow rate is proportional to the difference in stream function values.
\end{exercise}

\begin{exercise}
For the velocity field $\vec{v} = (-y, x)$, show that the streamlines are circles centered at the origin.
Compute the circulation around a circle of radius $R$ and verify that it equals $2\pi R^2$.
\end{exercise}

\begin{exercise}
Derive the velocity potential for flow past a cylinder and show that it satisfies Laplace's equation everywhere outside the cylinder.
\end{exercise}

\section{Stream Function Theory}

\subsection{Definition and Properties}

For two-dimensional incompressible flow, the \textbf{stream function} $\psi(x,y)$ is defined such that:
\[
u = \frac{\partial \psi}{\partial y}, \qquad v = -\frac{\partial \psi}{\partial x}
\]

This definition automatically satisfies the continuity equation $\frac{\partial u}{\partial x} + \frac{\partial v}{\partial y} = 0$.

\begin{theorem}[Stream Function Properties]
The following properties hold for any stream function $\psi$:
\\begin{enumerate}
    \item Streamlines are contours of constant $\psi$: $\psi(x,y) = C$.
    \item The volume flow rate between two streamlines equals $\psi_2 - \psi_1$.
    \item For irrotational flow, $\psi$ satisfies Laplace's equation: $\nabla^2 \psi = 0$.
    \item The velocity potential $\phi$ and stream function $\psi$ satisfy the Cauchy-Riemann equations.
\\end{enumerate}
\end{theorem}

\begin{lstlisting}[caption={Stream function for uniform flow}, label=lst:stream_uniform]
from euler import uniform_flow, stream_function

# Uniform flow at angle \alpha
vx, vy = uniform_flow(vx_val=1.0, vy_val=1.0)

# Stream function \psi = Uy - Vx for flow with components (U,V)
def psi_uniform(x, y):
    return y - x  # For vx=1, vy=1

print(f"\psi(1,2) = {psi_uniform(1, 2)}")
print(f"\psi(3,4) = {psi_uniform(3, 4)}")
\end{lstlisting}

\subsection{Stream Function for Elementary Flows}

\textbf{Uniform flow:} $\psi = Uy - Vx$\\
\textbf{Source at origin:} $\psi = \frac{m}{2\\pi}\\theta = \frac{m}{2\\pi}\\arctan(y/x)$\\
\textbf{Vortex at origin:} $\psi = -\frac{\\Gamma}{2\\pi}\\ln r = -\frac{\\Gamma}{4\\pi}\\ln(x^2+y^2)$\\
\textbf{Doublet:} $\psi = -\frac{\\kappa}{2\\pi}\frac{y}{x^2+y^2}$

These stream functions can be superposed to construct more complex flows.

\section{Differential Equations in Fluid Dynamics}

\subsection{The Euler Equations as ODEs}

Along a streamline, the Euler equations can be written as a system of ordinary differential equations. For steady, inviscid, incompressible flow:

\[
\frac{d\\vec{v}}{ds} = -\frac{1}{\\rho}\nabla p + \\vec{g}
\]

where $s$ is the arc length along the streamline.

\subsection{Vorticity Equation}

Taking the curl of Euler's equation gives the vorticity transport equation:

\[
\frac{D\\omega}{Dt} = (\\omega \\cdot \nabla)\\vec{v}
\]

For two-dimensional flow, this simplifies to $\frac{D\\omega}{Dt} = 0$, meaning vorticity is conserved following fluid particles.

\begin{lstlisting}[caption={Computing vorticity field}, label=lst:vorticity]
import numpy as np

def compute_vorticity(ux, uy, dx, dy):
    """Compute vorticity \omega = \partialv/\partialx - \partialu/\partialy from velocity fields"""
    dv_dx = np.gradient(uy, dx, axis=1)
    du_dy = np.gradient(ux, dy, axis=0)
    return dv_dx - du_dy

# Example: simple shear flow
x = np.linspace(0, 1, 100)
y = np.linspace(0, 1, 100)
X, Y = np.meshgrid(x, y)

# Shear flow: u = y, v = 0
ux = Y
uy = np.zeros_like(X)

omega = compute_vorticity(ux, uy, x[1]-x[0], y[1]-y[0])
print(f"Maximum vorticity: {np.max(np.abs(omega)):.4f}")
\end{lstlisting}

\subsection{Potential Flow Equations}

For irrotational flow ($\nabla \\times \\vec{v} = 0$), we can introduce a \textbf{velocity potential} $\phi$ such that $\\vec{v} = \nabla \phi$. Combined with incompressibility ($\nabla \\cdot \\vec{v} = 0$), this gives:

\[
\nabla^2 \phi = 0
\]

This is \textbf{Laplace's equation}, whose solutions are called \textbf{harmonic functions}. The stream function $\psi$ is also harmonic for irrotational flow.

\section{Boundary Value Problems}

\subsection{Flow Around Objects}

The classic problem of potential flow around a body requires solving Laplace's equation subject to boundary conditions:

\begin{itemize}
    \item \textbf{No-penetration:} $\frac{\partial \phi}{\partial n} = 0$ on the body surface
    \item \textbf{Far-field:} $\phi \to Ux$ as $r \to \infty$
\end{itemize}

\subsection{Numerical Solution}

Finite-difference methods can solve Laplace's equation on a grid using iterative techniques like Gauss-Seidel or successive over-relaxation (SOR).

\begin{lstlisting}[caption={Solving Laplace's equation}, label=lst:laplace]
import numpy as np

def solve_laplace(nx, ny, max_iter=1000, tol=1e-6):
    """Solve \nabla\textsuperscript{2}\varphi = 0 using Gauss-Seidel iteration"""
    phi = np.zeros((ny, nx))
    
    # Boundary conditions
    phi[-1, :] = 1.0  # Right: \varphi = Ux
    phi[0, :] = 0.0   # Left: \varphi = 0
    
    for iteration in range(max_iter):
        phi_old = phi.copy()
        
        for i in range(1, ny-1):
            for j in range(1, nx-1):
                phi[i, j] = 0.25 * (phi[i+1, j] + phi[i-1, j] + 
                                   phi[i, j+1] + phi[i, j-1])
        
        # Check convergence
        error = np.max(np.abs(phi - phi_old))
        if error < tol:
            print(f"Converged after {iteration} iterations")
            break
    
    return phi

phi = solve_laplace(50, 50)
print(f"\varphi at center: {phi[25, 25]:.4f}")
\end{lstlisting}

\section{Unsteady Flow and Differential Equations}

\subsection{Unsteady Streamlines}

For unsteady flow, streamlines are defined at each instant by solving:

\[
\frac{dx}{u(x,y,t)} = \frac{dy}{v(x,y,t)}
\]

This gives different streamline patterns at different times. The pathlines (trajectories of particles) are found by integrating:

\[
\frac{dx}{dt} = u(x,y,t), \qquad \frac{dy}{dt} = v(x,y,t)
\]

\subsection{Streaklines}

Streaklines connect all fluid particles that have passed through a fixed point. They are found by integrating the pathline equations backward in time from the release point.

\begin{lstlisting}[caption={Unsteady streamline computation}, label=lst:unsteady]
from euler import streamline_integrate

def unsteady_vx(x, y, t):
    return np.cos(t) * (1 - x*x + y*y)

def unsteady_vy(x, y, t):
    return np.sin(t) * (-2*x*y)

# Streamline at different times
for t in [0, np.pi/4, np.pi/2]:
    pts = streamline_integrate(
        lambda x, y: unsteady_vx(x, y, t),
        lambda x, y: unsteady_vy(x, y, t),
        0.0, 0.0, ds=0.02, n_steps=100
    )
    print(f"t = {t:.2f}: {len(pts)} points, end = ({pts[-1][0]:.2f}, {pts[-1][1]:.2f})")
\end{lstlisting}

\section{Applications}

\subsection{Aerodynamics}

Potential flow theory, despite its limitations, provides valuable insights for airfoil design. The \textbf{Kutta condition} selects the physically correct circulation around an airfoil, leading to the \textbf{Kutta-Joukowski theorem}:

\[
L' = \\rho U \\Gamma
\]

where $L'$ is the lift per unit span and $\\Gamma$ is the circulation.

\subsection{Meteorology}

Weather maps display streamlines of the atmospheric wind field. The geostrophic approximation relates wind to pressure gradients:

\[
f v = \frac{1}{\\rho}\frac{\partial p}{\partial x}, \qquad f u = -\frac{1}{\\rho}\frac{\partial p}{\partial y}
\]

where $f = 2\\Omega\\sin\phi$ is the Coriolis parameter.

\subsection{Biomedical Flow}

Blood flow in large arteries can be approximated as inviscid in the core region. Streamline analysis helps visualize flow patterns and identify regions of recirculation that may promote thrombus formation.

\section{Frequently Asked Questions}

\begin{itemize}
\item \textbf{Q: Why is the stream function useful?} \\\\
It automatically satisfies continuity for incompressible flow and reduces the number of dependent variables from two ($u,v$) to one ($\psi$). Contours of $\psi$ give streamlines directly.

\item \textbf{Q: Can stream functions be used for compressible flow?} \\\\
For compressible flow, a modified stream function can be defined using the mass flux $\\rho \\vec{v}$ instead of $\\vec{v}$. However, the resulting equations are more complex.

\item \textbf{Q: What is the difference between $\varphi$ and $\psi$?} \\\\
The velocity potential $\phi$ satisfies $\\vec{v} = \nabla\phi$ and exists for irrotational flow. The stream function $\psi$ satisfies $u = \psi_y, v = -\psi_x$ and exists for 2D incompressible flow. For irrotational, incompressible flow, both exist and are harmonic conjugates.

\item \textbf{Q: How do I handle singularities in stream functions?} \\\\
Singularities occur at sources, sinks, and vortices. In numerical simulations, these are typically avoided by placing them outside the computational domain or by using specialized grid treatments.

\item \textbf{Q: What is the relation between streamlines and equipotential lines?} \\\\
For irrotational, incompressible flow, streamlines ($\psi = \\text{const}$) and equipotential lines ($\phi = \\text{const}$) are orthogonal everywhere, forming a conformal grid.
\end{itemize}

\section{Advanced Exercises}

\begin{exercise}
Derive the stream function for a source-sink pair (dipole) and show that the streamlines are circles tangent to the $y$-axis at the origin.
\end{exercise}

\begin{exercise}
Show that for 2D irrotational flow, the stream function $\psi$ satisfies Laplace's equation $\nabla^2 \psi = 0$. Use this to derive the stream function for flow past a cylinder.
\end{exercise}

\begin{exercise}
Use the method of images to find the stream function for a source near a solid wall. Show that the wall is a streamline.
\end{exercise}

\begin{exercise}[Computational]
Implement a finite-difference solver for Laplace's equation on a grid with a circular cylinder. Compute the pressure distribution on the cylinder surface using Bernoulli's equation.
\end{exercise}

\begin{exercise}
For the velocity field $\\vec{v} = (x^2 - y^2, -2xy)$, find the velocity potential $\phi$ and stream function $\psi$. Verify that they satisfy the Cauchy-Riemann equations.
\end{exercise}

\begin{exercise}
Compute the stream function and velocity potential for flow past a rotating cylinder (with circulation $\\Gamma$). Show how the circulation affects the lift force via the Kutta-Joukowski theorem.
\end{exercise}

\begin{exercise}
Derive the vorticity equation from Euler's equations and show that for 2D flow, vorticity is materially conserved: $D\\omega/Dt = 0$.
\end{exercise}

\begin{exercise}
Use conformal mapping (Joukowski transformation) to map flow past a cylinder to flow past an airfoil. Describe the resulting flow pattern and lift characteristics.
\end{exercise}